% Options for packages loaded elsewhere
\PassOptionsToPackage{unicode}{hyperref}
\PassOptionsToPackage{hyphens}{url}
\documentclass[
]{article}
\usepackage{xcolor}
\usepackage{amsmath,amssymb}
\setcounter{secnumdepth}{-\maxdimen} % remove section numbering
\usepackage{iftex}
\ifPDFTeX
  \usepackage[T1]{fontenc}
  \usepackage[utf8]{inputenc}
  \usepackage{textcomp} % provide euro and other symbols
\else % if luatex or xetex
  \usepackage{unicode-math} % this also loads fontspec
  \defaultfontfeatures{Scale=MatchLowercase}
  \defaultfontfeatures[\rmfamily]{Ligatures=TeX,Scale=1}
\fi
\usepackage{lmodern}
\ifPDFTeX\else
  % xetex/luatex font selection
\fi
% Use upquote if available, for straight quotes in verbatim environments
\IfFileExists{upquote.sty}{\usepackage{upquote}}{}
\IfFileExists{microtype.sty}{% use microtype if available
  \usepackage[]{microtype}
  \UseMicrotypeSet[protrusion]{basicmath} % disable protrusion for tt fonts
}{}
\makeatletter
\@ifundefined{KOMAClassName}{% if non-KOMA class
  \IfFileExists{parskip.sty}{%
    \usepackage{parskip}
  }{% else
    \setlength{\parindent}{0pt}
    \setlength{\parskip}{6pt plus 2pt minus 1pt}}
}{% if KOMA class
  \KOMAoptions{parskip=half}}
\makeatother
\usepackage{color}
\usepackage{fancyvrb}
\newcommand{\VerbBar}{|}
\newcommand{\VERB}{\Verb[commandchars=\\\{\}]}
\DefineVerbatimEnvironment{Highlighting}{Verbatim}{commandchars=\\\{\}}
% Add ',fontsize=\small' for more characters per line
\newenvironment{Shaded}{}{}
\newcommand{\AlertTok}[1]{\textcolor[rgb]{1.00,0.00,0.00}{\textbf{#1}}}
\newcommand{\AnnotationTok}[1]{\textcolor[rgb]{0.38,0.63,0.69}{\textbf{\textit{#1}}}}
\newcommand{\AttributeTok}[1]{\textcolor[rgb]{0.49,0.56,0.16}{#1}}
\newcommand{\BaseNTok}[1]{\textcolor[rgb]{0.25,0.63,0.44}{#1}}
\newcommand{\BuiltInTok}[1]{\textcolor[rgb]{0.00,0.50,0.00}{#1}}
\newcommand{\CharTok}[1]{\textcolor[rgb]{0.25,0.44,0.63}{#1}}
\newcommand{\CommentTok}[1]{\textcolor[rgb]{0.38,0.63,0.69}{\textit{#1}}}
\newcommand{\CommentVarTok}[1]{\textcolor[rgb]{0.38,0.63,0.69}{\textbf{\textit{#1}}}}
\newcommand{\ConstantTok}[1]{\textcolor[rgb]{0.53,0.00,0.00}{#1}}
\newcommand{\ControlFlowTok}[1]{\textcolor[rgb]{0.00,0.44,0.13}{\textbf{#1}}}
\newcommand{\DataTypeTok}[1]{\textcolor[rgb]{0.56,0.13,0.00}{#1}}
\newcommand{\DecValTok}[1]{\textcolor[rgb]{0.25,0.63,0.44}{#1}}
\newcommand{\DocumentationTok}[1]{\textcolor[rgb]{0.73,0.13,0.13}{\textit{#1}}}
\newcommand{\ErrorTok}[1]{\textcolor[rgb]{1.00,0.00,0.00}{\textbf{#1}}}
\newcommand{\ExtensionTok}[1]{#1}
\newcommand{\FloatTok}[1]{\textcolor[rgb]{0.25,0.63,0.44}{#1}}
\newcommand{\FunctionTok}[1]{\textcolor[rgb]{0.02,0.16,0.49}{#1}}
\newcommand{\ImportTok}[1]{\textcolor[rgb]{0.00,0.50,0.00}{\textbf{#1}}}
\newcommand{\InformationTok}[1]{\textcolor[rgb]{0.38,0.63,0.69}{\textbf{\textit{#1}}}}
\newcommand{\KeywordTok}[1]{\textcolor[rgb]{0.00,0.44,0.13}{\textbf{#1}}}
\newcommand{\NormalTok}[1]{#1}
\newcommand{\OperatorTok}[1]{\textcolor[rgb]{0.40,0.40,0.40}{#1}}
\newcommand{\OtherTok}[1]{\textcolor[rgb]{0.00,0.44,0.13}{#1}}
\newcommand{\PreprocessorTok}[1]{\textcolor[rgb]{0.74,0.48,0.00}{#1}}
\newcommand{\RegionMarkerTok}[1]{#1}
\newcommand{\SpecialCharTok}[1]{\textcolor[rgb]{0.25,0.44,0.63}{#1}}
\newcommand{\SpecialStringTok}[1]{\textcolor[rgb]{0.73,0.40,0.53}{#1}}
\newcommand{\StringTok}[1]{\textcolor[rgb]{0.25,0.44,0.63}{#1}}
\newcommand{\VariableTok}[1]{\textcolor[rgb]{0.10,0.09,0.49}{#1}}
\newcommand{\VerbatimStringTok}[1]{\textcolor[rgb]{0.25,0.44,0.63}{#1}}
\newcommand{\WarningTok}[1]{\textcolor[rgb]{0.38,0.63,0.69}{\textbf{\textit{#1}}}}
\setlength{\emergencystretch}{3em} % prevent overfull lines
\providecommand{\tightlist}{%
  \setlength{\itemsep}{0pt}\setlength{\parskip}{0pt}}
\usepackage{bookmark}
\IfFileExists{xurl.sty}{\usepackage{xurl}}{} % add URL line breaks if available
\urlstyle{same}
% fallback for those not using the hyperref driver hyperxmp:
\makeatletter
\@ifundefined{xmpquote}{\newcommand{\xmpquote}[1]{#1}}{}
\makeatother
\hypersetup{
  hidelinks,
  pdfcreator={LaTeX via pandoc}}

\author{}
\date{}

\begin{document}

{
\setcounter{tocdepth}{2}
\tableofcontents
}
\section{WebCoder --- Contributing}\label{webcoder--contributing}

\begin{quote}
One-paragraph summary: WebCoder is a single-branch (\texttt{main})
monorepo --- Next.js 16 at the root, Django 6 under \texttt{backend/}
--- where every change lands through a PR that keeps BOTH test batteries
green (frontend: vitest + tsc + zero-warning lint + build + Playwright;
backend: pytest at 100 \% line coverage and warning-clean). This guide
explains the exact gates, the PR conventions, the BACKLOG workflow, and
how the roadmap "waves" operated.
\end{quote}

\textbf{Who it\textquotesingle s for:} anyone contributing code,
content, or docs. Romanian twin:
\href{../guide-ro/CONTRIBUIE.md}{\texttt{docs/guide-ro/CONTRIBUIE.md}}.

\begin{center}\rule{0.5\linewidth}{0.5pt}\end{center}

\subsection{1. Repository rules}\label{1-repository-rules}

\begin{itemize}
\tightlist
\item
  \textbf{\texttt{main} is the only base branch.} The old split
  (\texttt{backend} branch) was unified in October 2026; every PR
  targets \texttt{main}. Never reopen the split.
\item
  The repo is a monorepo: frontend at the root (\texttt{src/},
  \texttt{package.json}), backend under \texttt{backend/}, e2e under
  \texttt{e2e/}, mobile shell under \texttt{mobile/}.
\item
  Default merge is \textbf{squash}
  (\texttt{gh\ pr\ merge\ -\/-squash\ -\/-delete-branch}); true merge
  commits are reserved for history-preserving operations.
\item
  Work happens in disposable worktrees --- never commit inside a live
  deployment checkout.
\end{itemize}

\subsection{2. The gates (must pass before you
push)}\label{2-the-gates-must-pass-before-you-push}

Frontend --- run from the repo root:

\begin{Shaded}
\begin{Highlighting}[]
\ExtensionTok{npm}\NormalTok{ run test                      }\CommentTok{\# vitest}
\ExtensionTok{npx}\NormalTok{ tsc }\AttributeTok{{-}p}\NormalTok{ tsconfig.json }\AttributeTok{{-}{-}noEmit} \CommentTok{\# typecheck}
\ExtensionTok{npm}\NormalTok{ run lint                      }\CommentTok{\# 0 errors AND 0 warnings (zero{-}warnings policy)}
\ExtensionTok{npm}\NormalTok{ run build                     }\CommentTok{\# next build}
\ExtensionTok{npx}\NormalTok{ playwright test               }\CommentTok{\# e2e, mock{-}api mode (hermetic)}
\end{Highlighting}
\end{Shaded}

Backend --- any interpreter with \texttt{backend/deps.txt} installed:

\begin{Shaded}
\begin{Highlighting}[]
\VariableTok{SECRET\_KEY}\OperatorTok{=}\NormalTok{x }\VariableTok{DB\_PASSWORD}\OperatorTok{=}\NormalTok{x }\ExtensionTok{python} \AttributeTok{{-}m}\NormalTok{ pytest }\AttributeTok{{-}q}    \CommentTok{\# sqlite CI profile, 100 \% line cov, 0 warnings}
\end{Highlighting}
\end{Shaded}

Quality rules enforced on top:

\begin{itemize}
\tightlist
\item
  \textbf{Every NEW backend module reaches 100 \% line coverage} --- the
  repo-wide bar was 100 \% before; keep it.
\item
  \textbf{Zero warnings}: the pytest suite runs warning-clean (fix at
  the source, or add a narrowly-scoped, individually-commented
  \texttt{filterwarnings} ignore); \texttt{npm\ run\ lint} reports zero
  problems.
\item
  Ruff format + biome format run in CI; unformatted code fails the
  quality gate.
\item
  Migrations must not collide: waves number their migrations
  independently --- if two branches add the same-numbered migration,
  renumber the younger one before merge (see PRs \#360, \#375).
\end{itemize}

\subsection{3. PR conventions}\label{3-pr-conventions}

\begin{enumerate}
\def\labelenumi{\arabic{enumi}.}
\tightlist
\item
  Branch from \texttt{origin/main}
  (\texttt{git\ worktree\ add\ \textless{}path\textgreater{}\ -b\ \textless{}branch\textgreater{}\ origin/main}).
\item
  Implement \textbf{with tests} --- backend pytest for every new
  endpoint/command; frontend vitest for components; Playwright coverage
  for user-visible flows (extend \texttt{e2e/mock-api} fixtures; add
  real-stack specs when the feature touches the judge or auth).
\item
  Verify locally (the exact gates above), push, and open the PR against
  \texttt{main} with \texttt{gh\ pr\ create}.
\item
  Poll checks to green (CI runs \texttt{frontend.yml} on root-source
  paths, \texttt{django.yml} on \texttt{backend/**}, CodeQL on
  everything) --- up to 20 × 30 s.
\item
  \texttt{gh\ pr\ merge\ -\/-squash\ -\/-delete-branch}, then remove the
  worktree.
\item
  \textbf{Live-stack validation after backend merges} (RUNBOOK §PR
  conventions):
  \texttt{git\ -C\ \textless{}live\ worktree\textgreater{}\ pull\ -\/-ff-only\ origin\ main},
  rebuild \texttt{backend\ celery-worker\ judge-runner},
  \texttt{manage.py\ migrate}.
\end{enumerate}

Commit style: \texttt{type(scope):\ summary} --- e.g.
\texttt{feat(contests):\ …}, \texttt{fix(e2e):\ …},
\texttt{docs(guides):\ …}. The wave PRs are a good reference
(\texttt{git\ log\ -\/-oneline}).

\subsection{4. The BACKLOG workflow}\label{4-the-backlog-workflow}

\texttt{docs/BACKLOG.md} is the item-level tracker of record. Rules
(stated at the top of the file):

\begin{itemize}
\tightlist
\item
  statuses: \texttt{shipped\ (PR\ \#N)} · \texttt{in-flight\ (wave\ X)}
  · \texttt{planned\ (wave\ E/F/G)} · \texttt{someday};
\item
  \textbf{the wave agent that ships an item updates its status in the
  same PR};
\item
  numbering is stable --- never renumber; retire with
  \texttt{\textasciitilde{}}.
\end{itemize}

Above it sits \texttt{docs/SPEC-GAPS.md} (the platform backlog with
B/J/F/A/Q/SP/X ids) and \texttt{docs/ROADMAP-COMPETITIVE-2026-10-22.md}
(the 78-gap catalog + wave sequencing + the binding round-2
corrections). New ideas go to the BACKLOG parking lot; new platform gaps
get a SPEC-GAPS row.

\subsection{5. How the waves worked}\label{5-how-the-waves-worked}

The 2026-11 feature program ran as sequential waves (0 → A → B → C → D →
E → F → G), each a set of PRs implementing a roadmap slice:

\begin{enumerate}
\def\labelenumi{\arabic{enumi}.}
\tightlist
\item
  The wave\textquotesingle s spec was read from ROADMAP §4 + the round-2
  additions (ops-before-features, moderation-before-UGC, legal gates ---
  these override the roadmap when they disagree on order).
\item
  Items shipped as coherent PR groups (branch → tests → local verify →
  PR → checks → squash-merge), each PR updating BACKLOG statuses.
\item
  Judge/auth-touching items additionally validated on the live compose
  stack (RUNBOOK Mode A), including real-mode Playwright
  (\texttt{E2E\_API\_BASE=http://127.0.0.1:8010}).
\item
  Content work (the \textasciitilde1 850-problem import, hidden-test
  generation) ran as content-ops commands committed to the repo
  (\texttt{catalog\_polish}, \texttt{boost\_hidden\_tests},
  \texttt{export\_catalog}/\texttt{import\_catalog}) so every data
  mutation is reviewable and reproducible.
\end{enumerate}

\subsection{6. Content contributions}\label{6-content-contributions}

\begin{itemize}
\tightlist
\item
  Problems are authored \textbf{bilingually}: English statement
  mandatory (universal fallback), Romanian strongly preferred; 10
  further locales optional per problem.
\item
  Bac-style "write the subprogram" tasks are rewritten as whole-program
  I/O tasks when imported (same exercise, our own words --- see
  \texttt{problems/data/hidden\_test\_specs.py}).
\item
  Test suites: never import third-party test data you
  don\textquotesingle t have rights to; generate hidden suites via the
  committed spec format (reference solution + inputs + wrong-solution
  flip check).
\item
  The CNT-2 permission ledger (\texttt{ImportSource}) records source,
  terms and evidence for every imported corpus.
\end{itemize}

\subsection{7. Documentation
contributions}\label{7-documentation-contributions}

The docs suite is bilingual (EN under \texttt{docs/guide-en/}, RO under
\texttt{docs/guide-ro/}) and built into PDF/DOCX/TEX by
\texttt{docs/build-docs.ps1} (pandoc + tectonic; outputs committed under
\texttt{docs/dist/}). When you change behavior documented in a guide,
update both language twins in the same PR and rebuild the artifacts:

\begin{Shaded}
\begin{Highlighting}[]
\NormalTok{docs}\OperatorTok{/}\NormalTok{build{-}docs}\OperatorTok{.}\FunctionTok{ps1}    \CommentTok{\# regenerates docs/dist from the .md sources}
\end{Highlighting}
\end{Shaded}

Diagrams: Mermaid sources live in \texttt{docs/diagrams/*.mmd} with
rendered PNGs beside them --- edit the \texttt{.mmd}, re-render
(mermaid-cli), commit both.

\end{document}
